\documentclass[10pt]{article}
\usepackage[T1]{fontenc}
\usepackage[utf8]{inputenc}
\usepackage{lmodern,amsmath,amssymb,amsthm,booktabs,microtype}
\usepackage[margin=1in,includefoot]{geometry}
\usepackage[hidelinks]{hyperref}
\newtheorem{theorem}{Theorem}[section]
\newtheorem{lemma}[theorem]{Lemma}
\newtheorem{proposition}[theorem]{Proposition}
\newtheorem{corollary}[theorem]{Corollary}
\newcommand{\hsep}{\operatorname{hsep}}
\newcommand{\Req}{\operatorname{Req}}
\newcommand{\supp}{\operatorname{supp}}
\title{Four-Port Composition for Hamiltonian Edge Separation\\
and a Certificate-Based Induction Criterion}
\author{Thomas Bauer\\\small Barnette Search research note}
\date{Draft of September 17, 2026}
\begin{document}
\maketitle
\begin{abstract}
We describe the nine boundary states of a cubic fragment with four distinct
degree-two terminals. Compatible pairs of spanning path covers generate
Hamiltonian cycles across a matching four-edge cut. Selecting edge-union covers
within each state reduces the compatible product while preserving exactly its
set of covered ordered edge requirements. We also give a conditional
lift-and-repair criterion for an induction towards the conjectured extremal
bound. The boundary compatibility formalism is established in the literature;
priority is not asserted for the quantitative statements. The structural
existence and budget conditions needed for the extremal conjecture remain open.
\end{abstract}

\section{Fragments and the nine states}
A four-port fragment $X$ is a finite simple graph with four distinct labelled
terminals $x_1,x_2,x_3,x_4$, each of degree two, and every other vertex of degree
three. Let $Y$ be a vertex-disjoint fragment with terminals $y_1,\ldots,y_4$.
The indexing fixes the matching cut $T=\{t_i=x_i y_i:1\le i\le4\}$ of the
cubic graph $G=X\cup Y\cup T$. Write $I_X=E(X)$, $I_Y=E(Y)$ and
$n_X=|V(X)|$, $n_Y=|V(Y)|$; both orders are even. No planarity, bipartiteness
or connectivity premise is imposed on the fragments. Not every resulting
cubic graph is a Barnette graph.

An admissible path cover $F$ of $X$ is an acyclic spanning subgraph whose
degrees are two except at the terminals indexed by a set $S$, where they
are one. We require $|S|=2$ or $4$. Consequently $F$ consists of $|S|/2$
vertex-disjoint paths, has $n_X-|S|/2$ edges, and pairs the indices in $S$
according to the endpoints of these paths. Its state is $\alpha=(S,\pi)$.
There are nine possible formal states: six choices of a two-element $S$,
each with its unique pairing, and the three states
\[
 ([4],12|34),\qquad ([4],13|24),\qquad ([4],14|23).
\]
A fragment need not realize every state. In a disk embedding with terminals
in the order $1,2,3,4$, for example, the crossing pairing $13|24$ is impossible.

\begin{lemma}[Compatibility]\label{lem:compatibility}
Two covers with states $\alpha=(S,\pi)$ and $\beta=(S',\rho)$ glue to a
Hamiltonian cycle, using their selected cut edges, if and only if $S=S'$ and
either $|S|=2$, or $|S|=4$ and $\pi\ne\rho$.
Every Hamiltonian cycle of $G$ has a unique such representation.
\end{lemma}
\begin{proof}
Equal used-port sets are necessary for degree two at every terminal. With
two used ports, the two spanning paths join into one cycle. With four used
ports, contract each path to its terminal pairing and identify corresponding
terminals across the cut. Equal pairings give two components, hence two
cycles. Distinct perfect matchings on four indices form a single alternating
cycle, so the glued subgraph is connected and therefore Hamiltonian.

Conversely, a Hamiltonian cycle meets the nonempty proper cut $T$ an even,
positive number of times, hence twice or four times. Its restriction to
either side cannot contain a cycle component, which would disconnect it
from the other side. The restrictions are therefore admissible path covers.
Their edge sets and the used cut edges are recovered uniquely by restriction.
\end{proof}

Let $M_{\alpha\beta}$ denote this compatibility indicator. In the order with
the six two-port states first, its matrix is
\[
 M=\begin{pmatrix}I_6&0\\0&J_3-I_3\end{pmatrix}.
\]
Thus there are twelve compatible ordered state pairs. The matching
connectivity matrix is an established device for Hamiltonian-cycle
algorithms; see, for example, Curticapean, Lindzey and Nederlof~\cite{cln}.
This elementary four-port case is not claimed as a new connectivity method.

\section{A quantitative lemma preserving all covered requirements}
Choose finite families $\mathcal A_\alpha$ and $\mathcal B_\beta$ of admissible
covers in their indicated states, with no repetitions. Empty states are
allowed. Put $a_\alpha=|\mathcal A_\alpha|$ and $b_\beta=|\mathcal B_\beta|$.
For each nonempty state choose subfamilies $\mathcal U_\alpha$ and
$\mathcal V_\beta$ with the same edge union as their parent family:
\[
 \bigcup\mathcal U_\alpha=\bigcup\mathcal A_\alpha,\qquad
 \bigcup\mathcal V_\beta=\bigcup\mathcal B_\beta.
\]
Set $p_\alpha=|\mathcal U_\alpha|$ and $q_\beta=|\mathcal V_\beta|$;
the values for empty states are zero. These are covers of the edges that
actually occur in that state, not necessarily of every edge of the fragment.

For a family $\mathcal C$ of Hamiltonian cycles define
\[
 \Req(\mathcal C)=\{(e,f)\in E(G)^2:e\ne f,\ 
                    e\in C,\ f\notin C\text{ for some }C\in\mathcal C\}.
\]
The family strongly separates the edges precisely when this is the full set
of ordered distinct edge pairs. As in the main manuscript, $\hsep(G)=\infty$
when no such family exists.

\begin{theorem}[Requirement-preserving four-port composition]\label{thm:fourport}
Let $\mathcal K$ contain all gluings from compatible products
$\mathcal A_\alpha\times\mathcal B_\beta$. Let $\mathcal K'$ instead use,
in each compatible product, only
\[
 (\mathcal U_\alpha\times\mathcal B_\beta)
 \ \cup\ (\mathcal A_\alpha\times\mathcal V_\beta).
\]
Then
\begin{align}
 |\mathcal K|&=\sum_{\alpha,\beta}M_{\alpha\beta}a_\alpha b_\beta,
       \label{eq:full}\\
 |\mathcal K'|&=\sum_{\alpha,\beta}M_{\alpha\beta}
       (p_\alpha b_\beta+a_\alpha q_\beta-p_\alpha q_\beta)
       \le |\mathcal K|,\label{eq:refined}\\
 \Req(\mathcal K')&=\Req(\mathcal K).\label{eq:preservation}
\end{align}
For each nonempty state $\alpha=(S,\pi)$ the cover can be chosen with
\begin{equation}\label{eq:coverbound}
 1\le p_\alpha\le
 \min\{a_\alpha,n_X/2-1+|S|/2\},
\end{equation}
and similarly on $Y$. If $\mathcal K$ separates all edges, equation
\eqref{eq:refined} is an upper bound for $\hsep(G)$.
\end{theorem}
\begin{proof}
Lemma~\ref{lem:compatibility} gives Hamiltonicity and injectivity of gluing.
Within one product the intersection of the two selected rectangles is
$\mathcal U_\alpha\times\mathcal V_\beta$, proving both counts. The saving
in that block is $(a_\alpha-p_\alpha)(b_\beta-q_\beta)\ge0$.

We prove preservation one ordered pair at a time. Suppose $(e,f)$ has a
witness $C\star D\in\mathcal K$, with $C\in\mathcal A_\alpha$ and
$D\in\mathcal B_\beta$.
If both edges lie in $I_X\cup T$, retain $C$ and replace $D$ by any member
of $\mathcal V_\beta$. Such a member exists because the original block
is nonempty. Membership of both edges is preserved: the $X$ cover is
unchanged and cut-edge membership depends only on the state. The pair is
selected in $\mathcal A_\alpha\times\mathcal V_\beta$.
The symmetric argument treats pairs in $I_Y\cup T$. Together these cases
include every pair involving a cut edge.

If $e\in I_X$ and $f\in I_Y$, choose $C'\in\mathcal U_\alpha$ containing
$e$; it exists because the union agrees with that of $\mathcal A_\alpha$.
Keep $D$, which avoids $f$. The pair $(C',D)$ is selected and remains
compatible, so it witnesses $(e,f)$. For the reverse orientation, keep
the avoiding cover on $X$ and choose an edge-containing cover in
$\mathcal V_\beta$. Hence every previously covered requirement survives.
The reverse inclusion in \eqref{eq:preservation} follows from
$\mathcal K'\subseteq\mathcal K$.

Finally, $|I_X|=3n_X/2-2$. Start the union cover with one member of the
state, which has $n_X-|S|/2$ edges. At most
$n_X/2-2+|S|/2$ edges of the state's union are still uncovered. Add a member
containing a still-uncovered edge until none remain. Each step covers a new
edge and never repeats a member, proving \eqref{eq:coverbound}.
\end{proof}

\paragraph{Scope of the conclusion.}
The theorem preserves failures as well as successes: it does not establish
$\Req(\mathcal K)=\{(e,f):e\ne f\}$. If both input collections contain all
admissible covers, Lemma~\ref{lem:compatibility} identifies $\mathcal K$ with
the full Hamiltonian-cycle universe. Even then, finiteness of $\hsep$ is an
additional assertion. Nonempty compatible states alone do not suffice.
For instance, a family using all four cut edges in every cycle cannot
separate any two cut edges.

\section{An explicit criterion for separation}
The remaining existence condition can be stated using local data. For an
internal edge $e\in I_X$, define
\[
 A^+(e)=\{\alpha:\exists F\in\mathcal A_\alpha, e\in F\},\qquad
 A^-(e)=\{\alpha:\exists F\in\mathcal A_\alpha, e\notin F\},
\]
and define $B^+(f),B^-(f)$ similarly. Call a state active if it is nonempty
and has a compatible nonempty state on the other side. Represent a cover on
$X$ on the local ground set $I_X\cup\{t_1,t_2,t_3,t_4\}$ by adjoining
the formal cut edges indexed by its used-port set.

\begin{proposition}[Local and cross-edge criterion]\label{prop:criterion}
The full product $\mathcal K$ separates all edges if and only if:
\begin{enumerate}
\item The augmented covers in active states on $X$ separate every ordered
pair on $I_X\cup T$, and likewise for $Y$ on $I_Y\cup T$.
\item For every $e\in I_X$ and $f\in I_Y$, there is a compatible state pair
in $A^+(e)\times B^-(f)$ and a compatible state pair in
$A^-(e)\times B^+(f)$.
\end{enumerate}
\end{proposition}
\begin{proof}
Every active cover extends using any member of one compatible state. Thus
its local incidence information appears in $\mathcal K$, and every projected
global cycle is an active augmented cover. This proves equivalence for the
first set of requirements. For a cross-edge pair, independent choices of
an edge-containing cover and an edge-avoiding cover in compatible states
give exactly the witnesses described in the second condition. The cases
exhaust all ordered pairs.
\end{proof}

This criterion is checkable from local covers, but it is not asserted to
hold automatically for all Barnette fragments. It separates the structural
existence problem from the counting reduction in Theorem~\ref{thm:fourport}.

\section{Replacing a fragment: lift, repair and budget}
Let $H=X\cup Y_0\cup T_0$ and $G=X\cup Y_1\cup T_1$ share the same exterior
fragment $X$. Identify the four cut-edge labels, so that
$J=I_X\cup\{t_1,t_2,t_3,t_4\}$ is a common abstract edge set. Assume
$|V(H)|=m$, $|V(G)|=n>m$. Replacing a fragment need not preserve Barnette
properties; membership of both graphs in the chosen induction class must
be established separately.

\begin{lemma}[Lift and repair]\label{lem:lift}
Let $\mathcal C$ be a separating family of $H$, of size $k$. Suppose that
for each $C\in\mathcal C$, its exterior cover $C\cap I_X$ has a compatible
cover in $Y_1$, using the same cut-edge labels. Choose one such lift for
each $C$, and remove duplicates to obtain $\mathcal L$.
Let $\mathcal R$ be a family of $r$ Hamiltonian cycles of $G$ covering every
ordered requirement not covered by $\mathcal L$. Then
\[
 \hsep(G)\le|\mathcal L\cup\mathcal R|\le k+r.
\]
Every requirement needing repair involves at least one edge of $I_{Y_1}$.
\end{lemma}
\begin{proof}
The lifts exist by the stated compatibility hypothesis and are Hamiltonian
by Lemma~\ref{lem:compatibility}. For each old witness, all incidences on
$J$ remain unchanged. Hence all ordered pairs on $J$ remain separated,
even after duplicates are removed. The only remaining requirements have
an edge in $I_{Y_1}$. Adding their witnesses establishes separation, and
there are at most $k+r$ cycles.
\end{proof}

The repair family is a concrete upper certificate. It can be chosen from a
compatible product, or its reduced version in Theorem~\ref{thm:fourport},
provided that family covers the residual requirements. Neither existence
of a lift nor existence of the required repairs follows merely from the
number of available states. A greedy repair count is not an optimum.

Write $B(t)=\lfloor t(t+8)/32\rfloor$. If a known certificate in $H$ has
size $k\le B(m)$, its \emph{certified slack} is $s=B(m)-k\ge0$. This
quantity is computed from an explicit family, not from the unknown optimum.
Lemma~\ref{lem:lift} proves the desired bound whenever
\begin{equation}\label{eq:budget}
 r\le B(n)-B(m)+s.
\end{equation}
If this inequality is strict, then $\hsep(G)<B(n)$, which also excludes
$G$ from equality cases.

\begin{proposition}[Conditional induction criterion]\label{prop:induction}
Let $\mathcal G$ be a class of Barnette graphs whose orders are divisible
by four. For each $G\in\mathcal G$ let $\mathfrak A(G)$ specify the allowed
separating certificates, including any boundary-profile invariant required
by the reductions. Suppose:
\begin{enumerate}
\item There is a threshold $N_0$ such that every graph of the class with at
most $N_0$ vertices has a certificate in $\mathfrak A(G)$ of size at most
$B(|V(G)|)$.
\item Every larger $G\in\mathcal G$ has a smaller reduction
$H\in\mathcal G$ as above. For every certificate in $\mathfrak A(H)$ of
size $k\le B(m)$, a lift and repair as in Lemma~\ref{lem:lift} can be
chosen satisfying \eqref{eq:budget}, with resulting family in
$\mathfrak A(G)$.
\end{enumerate}
Then every $G\in\mathcal G$ has $\hsep(G)\le B(|V(G)|)$.
\end{proposition}
\begin{proof}
Induct on the order. The first condition supplies all base cases. The
induction hypothesis supplies an allowed certificate for the chosen smaller
graph. The second condition applies to that certificate; its output has
size at most $k+B(n)-B(m)+B(m)-k=B(n)$ and remains allowed.
\end{proof}

The universal quantifier over allowed certificates prevents a circular
argument that assumes a specially favorable smaller certificate without
showing that induction supplies one. An alternative is to track a stronger,
explicitly preserved boundary invariant. The proposition is elementary
bookkeeping; it is not a structural reduction theorem for Barnette graphs.
For the main manuscript's conjecture, one must still prove that an appropriate
class and reductions cover every order divisible by four from 40 onward.
Smaller orders cannot be silently used with this $B$: the cube already has
$\hsep=6>B(8)=4$, and the order-12 extremal value is $8>B(12)=7$.

\paragraph{Why a fixed increment is too restrictive.}
The proved double-ladder formula in the main manuscript gives
\[
 \hsep(D(3,15))=42,\qquad \hsep(D(5,15))=59.
\]
This square insertion increases the order from 40 to 44 and the exact value
by 17, although $B(44)-B(40)=71-60=11$. Thus a claim that \emph{every}
square insertion increases $\hsep$ by at most this difference is false.
The starting value leaves slack $60-42=18$, so no contradiction to the
extremal conjecture follows. These identities do not themselves construct
a lift with 17 repair cycles; that would require the separate lifting
and incidence checks in Lemma~\ref{lem:lift}.

\section{Independent checks and remaining research}
The bounded audit in \texttt{artifacts/four\_port\_20260917/} uses local
edge-subset enumeration to construct path covers. A separate standard-library
verifier checks the paths, pairings and edge-union covers, reconstructs global
cycles by actual connectivity, and independently enumerates the global
Hamiltonian universe through perfect-matching complements. It checks equality
of the full and reduced covered-requirement sets, including unsuccessful
separation cases. It imports neither the constructor nor the project solvers.

The audit covers all 24 terminal bijections for four pairs of fragments and
24 further examples with incomplete local families: 120 labelled cases,
not 120 isomorphism classes. Fragment orders are at most eight; resulting
graphs have at most 16 vertices. There are 16 Barnette cases, all separating,
and 16 strict reductions among all 120 cases. Four Barnette cases on 16
vertices have full-product size 14 and reduced upper certificates of size 13.
These bounds do not establish an exact value or the conjectured bound
$B(16)=12$. The fragments include a cube with two independent edges removed,
so the setup is not confined to pairs of ladders.

The regression suite additionally checks invalid pairings, empty union
covers, corrupted cycle certificates, failure to separate the cut when all
four cut edges are always used, and the fixed-increment obstruction. Commands,
versions, runtimes, labelled graph6 strings and SHA-256 hashes are recorded.
No random choices, optimization solver, or large census reconstruction is used.
Run from the repository root:
\begin{quote}\small\ttfamily
python artifacts/four\_port\_20260917/build\_examples.py\\
python -m pytest -q tests/test\_four\_port\_lemma.py
\end{quote}

The next substantive obligations are to find reductions preserving the
relevant graph class, prove lift availability and the local separation
criterion, and control repair costs with a preserved certificate invariant.
To establish uniqueness in Conjecture 11.2, a further analysis must show that
equality forces the balanced or nearly balanced double-ladder structure.
Neither the current state matrix nor the counting lemma supplies those
structural results. This note is an unreviewed proof draft, not a proof of
Conjecture 11.2 or Barnette's conjecture. AI tools assisted with derivation,
drafting and finite checks; the author remains responsible for the claims.

\begin{thebibliography}{9}
\bibitem{cln}
R. Curticapean, N. Lindzey, and J. Nederlof.
\emph{A Tight Lower Bound for Counting Hamiltonian Cycles via Matrix Rank}.
Author manuscript, 2017. \url{https://arxiv.org/abs/1709.02311}.
\bibitem{gsw}
M. Gorsky, R. Steiner, and S. Wiederrecht.
\emph{Matching theory and Barnette's conjecture}.
Discrete Mathematics 346(2), 113249, 2023.
\href{https://doi.org/10.1016/j.disc.2022.113249}{doi:10.1016/j.disc.2022.113249}.
The qualitative three-edge-cut context is developed in this work.
\end{thebibliography}
\end{document}
